\capitulo{6}{Trabajos relacionados}

Las herramientas de apoyo a los juegos de mesa pueden facilitar tanto el
aprendizaje inicial como la consulta durante la partida. Para ello recurren
a tutoriales, reglamentos enlazados o asistentes capaces de responder
preguntas, que ofrecen distintas formas de acceder a la información.

\section{Investigación relacionada}
\label{relacionados:investigacion}

Karim, Kang y Girouard estudian cómo facilitar el aprendizaje de juegos de
mesa a personas ciegas o con baja visión mediante un asistente de voz. A
partir de un primer estudio con catorce participantes, desarrollan
\textit{Board Game Assistant}, una aplicación de Alexa para
\textit{Ticket to Ride} que incorpora explicaciones breves, resúmenes y
opciones para adaptar la ayuda a las necesidades del jugador. Posteriormente,
nueve participantes valoran sus características mediante demostraciones
grabadas y materiales de referencia~\cite{karim2023rulebook}.

Entre sus conclusiones destaca la conveniencia de presentar las reglas
cuando hacen falta, permitir pausas y ofrecer información adicional a
petición del jugador, de modo que la explicación se adapte a su ritmo.
También señalan que la música y los recordatorios pueden distraer, por lo
que conviene poder ajustarlos. El estudio se limita a un juego y no evalúa
el uso directo del asistente durante partidas completas~\cite{karim2023rulebook}.

Su relación con \textit{Manualito} está en combinar una explicación inicial
con consultas puntuales, adaptando la cantidad de información a lo que el
jugador necesita en cada momento.

\section{Aplicaciones relacionadas}
\label{relacionados:aplicaciones}

\subsection{Dized}
\label{relacionados:dized}

\textit{Dized} acompaña a los jugadores mediante tutoriales que combinan
animaciones, voz y texto para presentar las reglas de forma progresiva
mientras juegan. Además, ofrece reglamentos digitales con buscador,
menús organizados y preguntas frecuentes, que pueden incorporar
ampliaciones y correcciones~\cite{dized_funciones}.

Sus herramientas permiten preparar contenidos específicos para cada juego,
ya sea desde la editorial o con ayuda de otros
creadores~\cite{dized_funciones}. Esta preparación
facilita una enseñanza guiada, mientras que \textit{Manualito} genera
explicaciones a partir de los manuales incorporados por los usuarios, que
pueden ampliarlas mediante preguntas.

\subsection{Rulepop}
\label{relacionados:rulepop}

\textit{Rulepop} transforma el reglamento de cada juego en una web que
combina resúmenes, referencias de iconos y reglas completas enlazadas entre
sí. Así, el jugador puede partir de una explicación breve y acceder al
detalle correspondiente, o seguir un enlace hacia otra regla relacionada.
Cada sitio puede instalarse como aplicación y consultarse sin
conexión~\cite{rulepop_funciones}.

La preparación incluye reorganizar el contenido, revisar errores e
incorporar las observaciones de la editorial~\cite{rulepop_funciones}.
Su propuesta muestra cómo
facilitar la consulta mediante la estructura del reglamento y la navegación
entre sus entradas, mientras que \textit{Manualito} permite expresar una
duda con palabras propias y obtener una respuesta acompañada de las
páginas utilizadas.

\subsection{RulesBot.ai}
\label{relacionados:rulesbot}

\textit{RulesBot.ai} ofrece un asistente al que plantear preguntas sobre
juegos de mesa y describe sus respuestas como explicaciones acompañadas
de citas~\cite{rulesbot_presentacion}. Su catálogo da acceso a un chat
asociado a cada juego, por lo que la consulta comienza seleccionando el
título correspondiente~\cite{rulesbot_catalogo}.

Comparte con \textit{Manualito} la consulta conversacional por juego,
aunque la documentación pública no aclara si permite subir manuales,
corregir su texto o conservar conversaciones. Por ello, estas funciones
quedan sin verificar y no se consideran ausentes.

\section{Comparación con \textit{Manualito}}
\label{relacionados:comparacion}

La tabla~\ref{tabla:relacionados-comparacion} resume las características
descritas en la documentación pública de las aplicaciones, consultada el
11 de septiembre de 2026, y las implementadas en \textit{Manualito}.
La comparación se centra en las posibilidades de uso, ya que valorar la
rapidez o la precisión de las respuestas requeriría una evaluación común.

\begin{table}[htbp]
    \centering
    \small
    \begin{tblr}{
        width=\textwidth,
        colspec={X[1.2,l] X[1,l] X[1,l] X[1.1,l] X[1.3,l]},
        cells={valign=t},
        row{1}={font=\bfseries},
        column{1}={font=\bfseries},
        colsep=4pt,
        rowsep=4pt
    }
        \toprule
        Aspecto & Dized & Rulepop & RulesBot.ai & Manualito \\
        \midrule
        Ayuda inicial & Tutoriales progresivos & Resúmenes enlazados & Explicaciones mediante chat & Explicación general del juego \\
        Consulta de dudas & Buscador, menús y preguntas frecuentes & Navegación entre reglas & Preguntas mediante chat & Preguntas y conversación posterior \\
        Contenido & Preparado para cada juego & Preparado con la editorial & Juegos del catálogo & Manuales aportados por usuarios \\
        Acceso a las fuentes & Reglas digitales & Entradas enlazadas & Citas anunciadas & Referencias a las páginas originales \\
        Organización & Catálogo de juegos & Un sitio por juego & Chat por juego & Juegos, manuales y conversaciones \\
        \bottomrule
    \end{tblr}
    \caption{Comparación de las aplicaciones relacionadas con \textit{Manualito}.
    Elaboración propia a partir de la documentación de
    Dized~\cite{dized_funciones}, Rulepop~\cite{rulepop_funciones} y
    RulesBot.ai~\cite{rulesbot_presentacion,rulesbot_catalogo}.}
    \label{tabla:relacionados-comparacion}
\end{table}
\FloatBarrier

\textit{Manualito} reúne la explicación inicial, la consulta conversacional
y la gestión de los manuales en un mismo espacio. El usuario puede añadir
documentos en \acs{PDF} o imágenes, revisar el texto obtenido y corregirlo,
mientras que los juegos y las conversaciones quedan guardados para volver
a ellos en otra partida. Esto permite incorporar un manual sin tener que
elaborar un tutorial específico, aunque primero hay que cargarlo y esperar
a que termine su procesamiento. Además, la explicación generada no ofrece
el acompañamiento paso a paso de un tutorial preparado como los de Dized.
